\section{第 25 轮：五项亏格恒等式、偶极子的余弦形式与全局窗口归约}\label{sec:r25}

前二十四轮的绝大部分工作都是\emph{分截面}进行的：先固定 $\Lam$ 的谱型（中心截面 $\Lam=I_3/3$、
单轴截面 $\Lam=bI_3+\rho ee^\top$、秩一截面 $\Lam=zz^\top$），再在该截面上做矩层级或分块分解。
第 24 轮把中心截面彻底闭合（定理~\ref{thm:r24main}），但同时暴露出一个结构性障碍
（注~\ref{rem:r24obs}）：单轴截面上「目标权重 $w_i=g_ih_i$」与「各向同性权重 $g_i$」分离，
第 24 轮的全部分块恒等式都建立在 $\sum_im_iu_i=1$ 这一等式上，而它在倾斜权重下不再成立。

本轮换一条完全不同的路。我们不再区分截面，而是对\textbf{任意}三元组 $(g,U,\Lam)$
（$(g,U)$ 满足 \eqref{eq:frame18}，$\Lam\succeq0$、$\tr\Lam=1$）建立一个\emph{恒等式}，
把 $2\rho_0-J_w$ 写成「一个由五项非负量组成的亏格 $\Phi$」减去常数 $\delta_0/3$：
\begin{equation}\label{eq:r25head}
\boxed{\;2\rho_0-J_w=\Phi-\frac{\delta_0}{3}\;},\qquad
J_w:=L_w+A_w^2 .
\end{equation}
于是全局猜想 $\kappa_3\le\rho_0$ 被\textbf{精确等价}地改写成一条不含图、不含阶数、
不含掩码、也不含截面分类的单项不等式 $\Phi\ge\delta_0/3$（定理~\ref{thm:r25main}、
推论~\ref{cor:r25sharp}）。这个改写有三点价值。
第一，$\Phi$ 的五项\emph{逐项非负}，因此可以逐项丢弃得到无条件结论：
丢掉四项就得到\textbf{窗口定理}（定理~\ref{thm:r25window}）——
只要 $h$ 的色散足够大（$\sin^2\theta_1\ge\delta_0/(9c_W)$，即 $W\le W_c=1.666927532150$），
$\kappa_3\le\rho_0$ \emph{无条件成立}。全局问题由此被压缩到
$\theta_1<15.7614^\circ$ 的一条薄窗口内，而窗口的左端点 $\theta_1=0$ 正是第 24 轮已闭合的中心截面。
第二，冠军构型在 \eqref{eq:r25head}~中\textbf{精确取等}（$\Phi=\delta_0/3$，差 $6.9\times10^{-17}$），
且五项\emph{全部严格为正}（表~\ref{tab:r25budget}）。这给出一个非常强的\emph{否定性}信息：
任何在冠军处不取等的松弛都注定失败，任何「先丢掉一项再估计其余四项」的做法都不可能证出全局结论。
第三，$\Phi$ 的第一项 $\sin^2\theta=1-A_w^2$ 有一个干净的几何解释（引理~\ref{lem:r25dipole}）：
$A_w=1-E/2$，其中 $E$ 是 $h$ 到「秩一剖面」集合 $\{a_z:z\in S^2\}$ 在 $L^2(\nu)$ 中的距离平方，
于是 $A_w=\cos\theta$ 是 $h$ 与最近秩一剖面的 $L^2(\nu)$-夹角余弦，并且\textbf{$A_w\le1$ 无条件成立}。
这条不等式本身就是本轮的一个独立结果：它说明「偶极子」通道
$\sum_iw_i|\langle u_i,z\rangle|$ 永远不超过总质量 $1$，等号恰在 $\Lam$ 秩一时成立（命题~\ref{prop:r25rank1}）。

\subsection{记号与 $L^2(\nu)$ 视角}\label{sec:r25:not}

设 $N\ge3$，$(g,U)$ 满足 \eqref{eq:frame18}，$\Lam\succeq0$、$\tr\Lam=1$。
把帧测度 $\nu:=\sum_{i=1}^Ng_i\delta_{u_i}$ 看作 $S^2$ 上的正测度（总质量 $\sum_ig_i=3$），
并引入加权内积
\[
\langle f,f'\rangle_g:=\int_{S^2}ff'\dd\nu=\sum_ig_if(u_i)f'(u_i),
\qquad \|f\|_g^2=\langle f,f\rangle_g .
\]
对本轮而言，$L^2(\nu)$ 中有两类特殊函数：
\textbf{坐标函数} $a_k:=\langle\,\cdot\,,e_k\rangle$（$k=1,2,3$，$e_k$ 为 $\Lam$ 的特征标架），
以及\textbf{秩一剖面} $a_z:=|\langle\,\cdot\,,z\rangle|$（$z\in S^2$）。
定义
\begin{equation}\label{eq:r25def}
\begin{aligned}
&h:=\sqrt{\textstyle\sum_k\lambda_ka_k^2}\ \ (\text{即 }h_i=\sqrt{u_i^\top\Lam u_i}),\qquad
w_i:=g_ih_i,\qquad W:=\sum_iw_i,\\
&S_w:=\sum_iw_iu_iu_i^\top,\qquad
L_w:=\sum_{ij}w_iw_j|\langle u_i,u_j\rangle|,\qquad
A_w:=\max_{\|z\|=1}\sum_iw_i|\langle u_i,z\rangle|,\\
&M_4:=\sum_{ij}w_iw_j\langle u_i,u_j\rangle^4,\qquad
Z_k:=\sum_{ij}w_iw_jP_k(\langle u_i,u_j\rangle),\\
&D_w:=\sum_{ij}w_iw_j\bigl[q(\langle u_i,u_j\rangle^2)-|\langle u_i,u_j\rangle|\bigr],\\
&H:=\sum_{ij}g_ig_j(h_i-h_j)^2,\qquad
\widetilde E:=\sum_{ij}g_ig_j(h_i-h_j)^2\langle u_i,u_j\rangle^2,\\
&E(z):=\sum_ig_i\bigl(h_i-a_i(z)\bigr)^2,\qquad E:=\min_{\|z\|=1}E(z),
\end{aligned}
\end{equation}
其中 $P_k$ 是三维勒让德多项式（$P_2(t)=\frac12(3t^2-1)$、$P_4(t)=\frac18(35t^4-30t^2+3)$），
$q(s)=\gamma+\alpha s+\beta s^2$ 是第 21 轮的中心主函数
（$\gamma=\frac{1+\sqrt5}{16}$、$\alpha=\frac{7\sqrt5-5}8$、$\beta=\frac{25-15\sqrt5}{16}<0$，
满足 $q(s)\ge\sqrt s$ 于 $[0,1]$，故 $D_w\ge0$）。
沿用命题~\ref{prop:r24gen}，$\rho_{g,U}(\Lam)=\frac12[L_w+\alpha_w]$，
且对上确界而言 $\alpha_w$ 可换成符号最优值 $A_w^2$，故
\begin{equation}\label{eq:r25Jw}
\kappa_3=\tfrac12\sup_{(g,U,\Lam)}J_w,\qquad J_w:=L_w+A_w^2 .
\end{equation}
最后引入三个常数与两个角：
\begin{equation}\label{eq:r25const}
c_W:=\gamma-\frac{3\beta}{35}=\frac{2\sqrt5-1}{14},\quad
c_S:=\alpha+\frac{6\beta}7=\frac{10+\sqrt5}{14},\quad
|c_Z|:=-\frac{8\beta}{35}=\frac{3\sqrt5-5}{14},
\end{equation}
\begin{equation}\label{eq:r25ang}
\sin^2\theta:=1-A_w^2,\qquad
\cos\theta_1:=\frac{W}{\sqrt3},\quad \sin^2\theta_1=\frac{3-W^2}3=\frac H6 .
\end{equation}
数值上 $c_W=0.2480097111$、$c_S=0.8740048555$、$|c_Z|=0.1220145666$，
并且下列三条精确代数关系是本轮全部计算的枢纽（$60$ 位核对）：
\begin{equation}\label{eq:r25alg}
3c_W+c_S=\varphi,\qquad
\varphi+1-2\rho_0=\frac{\delta_0}3,\qquad
2\rho_0-1=\varphi-\frac{\delta_0}3 ,
\end{equation}
其中 $\delta_0=3\varphi+3-6\rho_0=\frac{9+7\sqrt5-4\sqrt{35}}6=0.164692785017$，
$\frac{\delta_0}3=0.054897595005559$。

\subsection{四个预备引理}\label{sec:r25:lem}

\begin{lemma}[正交坐标基]\label{lem:r25ortho}
在 \eqref{eq:frame18}~下，坐标函数 $a_1,a_2,a_3$ 在 $L^2(\nu)$ 中标准正交：
$\langle a_k,a_l\rangle_g=\delta_{kl}$。因此
\begin{equation}\label{eq:r25ortho1}
h^2=\sum_k\lambda_ka_k^2,\qquad \|h\|_g^2=\tr\Lam=1,\qquad
\|a_z\|_g^2=z^\top\Bigl(\sum_ig_iu_iu_i^\top\Bigr)z=1\ \ \forall z\in S^2 .
\end{equation}
\end{lemma}

\begin{proof}
$\langle a_k,a_l\rangle_g=\sum_ig_i(u_i)_k(u_i)_l=(\sum_ig_iu_iu_i^\top)_{kl}=\delta_{kl}$。
$h^2(u)=u^\top\Lam u=\sum_k\lambda_k\langle u,e_k\rangle^2$。
$\|h\|_g^2=\sum_ig_iu_i^\top\Lam u_i=\tr(\Lam\sum_ig_iu_iu_i^\top)=\tr\Lam$。
\end{proof}

引理~\ref{lem:r25ortho}~的内容是平凡的，但它把问题放进了一个 Hilbert 空间：
$h$ 与所有 $a_z$ 都是\textbf{单位向量}。这正是下一条引理能够成立的原因。

\begin{lemma}[偶极子的余弦形式，本轮核心]\label{lem:r25dipole}
对每个 $z\in S^2$，
\begin{equation}\label{eq:r25dip1}
\sum_iw_ia_i(z)=1-\frac12E(z)=1-\frac12\sum_ig_i\bigl(h_i-a_i(z)\bigr)^2 .
\end{equation}
从而
\begin{equation}\label{eq:r25dip2}
A_w=1-\frac E2\le1,\qquad
\sin^2\theta=1-A_w^2=E-\frac{E^2}4,\qquad
A_w\ge\Bigl\|\sum_iw_iu_i\Bigr\| ,
\end{equation}
且 $A_w=\cos\theta$，$\theta$ 是 $h$ 与秩一剖面族 $\{a_z\}$ 在 $L^2(\nu)$ 中的最小夹角。
若最大值在胞腔内部取得（$z_*$ 使 $m:=\sum_iw_i\varepsilon_i(z_*)u_i$ 与 $z_*$ 平行），则
\begin{equation}\label{eq:r25dip3}
A_w^2=W\zeta-\frac12V_w,\qquad
\zeta:=z_*^\top S_wz_*,\quad V_w:=\sum_{ij}w_iw_j\bigl(a_i(z_*)-a_j(z_*)\bigr)^2 .
\end{equation}
\end{lemma}

\begin{proof}
展开平方并用引理~\ref{lem:r25ortho}：
\[
E(z)=\|h\|_g^2-2\langle h,a_z\rangle_g+\|a_z\|_g^2
=1-2\sum_ig_ih_i|\langle u_i,z\rangle|+1
=2-2\sum_iw_ia_i(z),
\]
即 \eqref{eq:r25dip1}。因 $E(z)\ge0$ 且 $A_w=\max_z\sum_iw_ia_i(z)=1-\frac12\min_zE(z)=1-\frac E2$，
得 $A_w\le1$；代入 $1-A_w^2=1-(1-\frac E2)^2=E-\frac{E^2}4$。
由 $\|h\|_g=\|a_z\|_g=1$，$\langle h,a_z\rangle_g=\sum_iw_ia_i(z)$ 正是余弦，
故 $\theta$ 的解释成立。
最后一句：对任意 $z$，$A_w\ge\sum_iw_i|\langle u_i,z\rangle|\ge|\langle m,z\rangle|$，
取 $z=m/\|m\|$ 得 $A_w\ge\|m\|$。
\eqref{eq:r25dip3}：$V_w=2(\sum_iw_ia_i^2)W-2(\sum_iw_ia_i)^2=2W\zeta-2A_w^2$，
其中 $\sum_iw_ia_i(z_*)^2=z_*^\top S_wz_*=\zeta$。
\end{proof}

\begin{remark}[引理~\ref{lem:r25dipole}~的含义]\label{rem:r25dip}
\eqref{eq:r25dip1}~是平行四边形恒等式在 $L^2(\nu)$ 中的直接后果，
但它的内容并不平凡：它把「偶极子通道的最大值 $A_w$」与「$\Lam$ 偏离秩一的程度 $E$」
用\textbf{等式}绑在一起。三点推论值得记下。
(a) $A_w\le1$ 是一条\emph{无条件}的通道上界，比 \STW\ 框架里逐点估计所得的
$A_w\le\sum_iw_i\le\sqrt3$ 强得多（$\sqrt3=1.732$ 对 $1$）；
(b) $E=0\iff h=a_z\iff\Lam=zz^\top$（秩一），故秩一截面是\emph{刚性}的：
$A_w\equiv1$ 对一切几何成立（命题~\ref{prop:r25rank1}）；
(c) $A_w$ 与 $W$ 是两个\emph{独立}的通道量：$W=\sum w_i$ 只涉及 $h$ 的一阶矩，
$A_w$ 涉及 $h$ 到秩一剖面的距离。第 24 轮之所以只能处理中心截面，
本质上是因为那里 $h\equiv\frac1{\sqrt3}$ 使二者退化；本轮把它们分开对待。
\end{remark}

\begin{lemma}[通道亏格恒等式]\label{lem:r25chan}
下列四个关系对一切 $(g,U,\Lam)$ \textbf{精确}成立：
\begin{align}
W^2&=3-\frac H2\qquad(\text{即 }H=6-2W^2=6\sin^2\theta_1),\label{eq:r25ch1}\\
\|S_w\|_F^2&=1-\frac{\widetilde E}2,\label{eq:r25ch2}\\
\|S_w\|_F^2&\ge\frac{W^2}3\iff\widetilde E\le\frac H3
\qquad(\text{二者等价，均为 }Z_2\ge0\text{ 的推论}),\label{eq:r25ch3}\\
A_w^2&=W\zeta-\frac12V_w\quad(\text{引理~\ref{lem:r25dipole}}).\label{eq:r25ch4}
\end{align}
\end{lemma}

\begin{proof}
\eqref{eq:r25ch1}：$H=\sum_{ij}g_ig_j(h_i-h_j)^2=2(\sum_ig_i)(\sum_ig_ih_i^2)-2(\sum_ig_ih_i)^2$，
而 $\sum_ig_i=\tr I_3=3$、$\sum_ig_ih_i^2=\tr\Lam=1$、$\sum_ig_ih_i=W$，故 $H=6-2W^2$。
\eqref{eq:r25ch2}：把 $\|S_w\|_F^2=\sum_{ij}w_iw_j\langle u_i,u_j\rangle^2$ 与 $\widetilde E$ 相加，
\[
\|S_w\|_F^2+\frac{\widetilde E}2
=\sum_{ij}g_ig_j\langle u_i,u_j\rangle^2\Bigl[h_ih_j+\frac12(h_i-h_j)^2\Bigr]
=\frac12\sum_{ij}g_ig_j\langle u_i,u_j\rangle^2(h_i^2+h_j^2),
\]
因为 $h_ih_j+\frac12(h_i-h_j)^2=\frac12(h_i^2+h_j^2)$。上式按 $i$ 求和先做：
$\sum_jg_j\langle u_i,u_j\rangle^2=u_i^\top(\sum_jg_ju_ju_j^\top)u_i=u_i^\top I_3u_i=1$，
故右端 $=\frac12\sum_ig_ih_i^2\cdot2\cdot\frac12\cdot2=\sum_ig_ih_i^2=\tr\Lam=1$。
\eqref{eq:r25ch3}：核 $P_2$ 在 $S^2$ 上正定（Schoenberg），故
$Z_2=\sum_{ij}w_iw_jP_2(\langle u_i,u_j\rangle)\ge0$；
而 $Z_2=\frac12(3\|S_w\|_F^2-W^2)$。
另一方面 $\widetilde E-\frac H3=2-2\|S_w\|_F^2-2+\frac{2W^2}3=-2(\|S_w\|_F^2-\frac{W^2}3)$，故二者等价。
\end{proof}

\begin{remark}
\eqref{eq:r25ch2}~是本轮技术上最关键的一步：它说明
$\|S_w\|_F^2$（一个纯角向量）与 $\widetilde E$（一个含 $h$-色散的量）之和\textbf{恒等于 $1$}。
这个恒等式的来源是各向同性条件被用了\emph{两次}：
一次在 $\sum_jg_j\langle u_i,u_j\rangle^2=1$，一次在 $\sum_ig_ih_i^2=\tr\Lam$。
它把第 24 轮注~\ref{rem:r24obs}~指出的「两套权重分离」的困难\emph{绕过}了：
我们不需要 $\sum_iw_iu_iu_i^\top=I_3$（这是错的），
只需要 $\sum_i(w_i/h_i)u_iu_i^\top=I_3$（这是命题~\ref{prop:r24gen}\eqref{eq:r24gen3}~给出的），
而它恰好足以产生 \eqref{eq:r25ch2}。
\end{remark}

\begin{lemma}[Schoenberg 四阶分解与关键消去]\label{lem:r25scho}
在 $S^2$ 上有精确展开
\begin{equation}\label{eq:r25scho1}
t^4=\frac15P_0(t)+\frac47P_2(t)+\frac8{35}P_4(t),
\end{equation}
故
\begin{equation}\label{eq:r25scho2}
M_4=-\frac3{35}W^2+\frac67\|S_w\|_F^2+\frac8{35}Z_4,\qquad Z_4\ge0 ,
\end{equation}
并且
\begin{equation}\label{eq:r25scho3}
L_w=\gamma W^2+\alpha\|S_w\|_F^2+\beta M_4-D_w ,
\end{equation}
\begin{equation}\label{eq:r25key}
\boxed{\;D_w+|c_Z|Z_4=c_WW^2+c_S\|S_w\|_F^2-L_w\;}
\end{equation}
其中 $c_W,c_S,|c_Z|$ 如 \eqref{eq:r25const}，三者\textbf{全部为正}。
\end{lemma}

\begin{proof}
\eqref{eq:r25scho1}：$\frac47P_2=\frac{6t^2-2}7$，$\frac8{35}P_4=t^4-\frac67t^2+\frac3{35}$，
相加得 $t^4+\frac15-\frac27+\frac3{35}=t^4$。
代入 $\sum_{ij}w_iw_j(\cdot)$ 并用 $\sum_{ij}w_iw_jP_2=Z_2=\frac12(3\|S_w\|_F^2-W^2)$ 得
$M_4=\frac15W^2+\frac47Z_2+\frac8{35}Z_4=-\frac3{35}W^2+\frac67\|S_w\|_F^2+\frac8{35}Z_4$；
$Z_4\ge0$ 仍是 Schoenberg（$P_4$ 在 $S^2$ 上正定）。
\eqref{eq:r25scho3}：$\sum_{ij}w_iw_jq(\langle u_i,u_j\rangle^2)
=\gamma W^2+\alpha\sum_{ij}w_iw_j\langle u_i,u_j\rangle^2+\beta M_4
=\gamma W^2+\alpha\|S_w\|_F^2+\beta M_4$，
而 $D_w$ 的定义正是它与 $L_w$ 之差。
\eqref{eq:r25key}：把 \eqref{eq:r25scho2}~代入 \eqref{eq:r25scho3}，
$\beta\cdot\frac8{35}Z_4=-\frac{8|\beta|}{35}Z_4=-|c_Z|Z_4$（因 $\beta<0$），
$\beta\cdot(-\frac3{35})=-\frac{3\beta}{35}$，$\beta\cdot\frac67=\frac{6\beta}7$，于是
\[
L_w=\Bigl(\gamma-\frac{3\beta}{35}\Bigr)W^2+\Bigl(\alpha+\frac{6\beta}7\Bigr)\|S_w\|_F^2
-|c_Z|Z_4-D_w=c_WW^2+c_S\|S_w\|_F^2-|c_Z|Z_4-D_w .
\]
移项即 \eqref{eq:r25key}。三个常数的正性：
$c_W=\frac{2\sqrt5-1}8\cdot\frac8{14}>0$（$2\sqrt5>1$），
$c_S>0$ 显然，$|c_Z|=\frac{3\sqrt5-5}{14}>0$（$45>25$）。
\end{proof}

\begin{remark}[为什么 \eqref{eq:r25key}~是本轮的枢纽]\label{rem:r25key}
主函数 $q$ 的二次项系数 $\beta$ 是\textbf{负}的（$\beta=-0.533813$），
这是第 21--24 轮所有估计中最难处理的一项：
$L_w=\gamma W^2+\alpha\|S_w\|_F^2+\beta M_4-D_w$ 中 $\beta M_4$ 是负的，
必须用 $M_4$ 的\emph{上}界来控制它，而 $M_4$ 的上界又依赖于 $Z_4\ge0$ 之外的信息。
\eqref{eq:r25scho1}~的分解说明：$M_4$ 中唯一「不受 $W^2$ 与 $\|S_w\|_F^2$ 控制」的部分
恰好是 $\frac8{35}Z_4$，而它带的系数 $\frac{8\beta}{35}$ 是负的，
于是 $-\frac{8|\beta|}{35}Z_4$ 与 $-D_w$ 一起\textbf{移到左边变成两个非负项}。
换句话说，$\beta<0$ 造成的全部损失被精确地兑换成了
$|c_Z|Z_4+D_w\ge0$，而剩下的系数 $c_W,c_S$ 反而\emph{变正}了。
这不是放缩，是恒等式；这正是它比第 22--24 轮的矩松弛更有力的原因。
\end{remark}

\subsection{主恒等式与全局归约}\label{sec:r25:main}

\begin{definition}[五项亏格]\label{def:r25phi}
\begin{equation}\label{eq:r25phi}
\Phi:=\underbrace{\sin^2\theta}_{\text{秩一亏格}}
+\underbrace{3c_W\sin^2\theta_1}_{=\,(c_W/2)H}
+\underbrace{\frac{c_S}2\widetilde E}_{\text{角向稀释}}
+\underbrace{|c_Z|Z_4}_{\text{四阶 Schoenberg}}
+\underbrace{D_w}_{\text{主函数亏格}} .
\end{equation}
五项均 $\ge0$：第一项由引理~\ref{lem:r25dipole}~的 $A_w\le1$，
第二、三项由 $H,\widetilde E\ge0$ 与 $c_W,c_S>0$，
第四项由 Schoenberg，第五项由 $q(s)\ge\sqrt s$。
\end{definition}

\begin{theorem}[五项亏格主恒等式与全局归约，本轮主定理]\label{thm:r25main}
对一切 $N\ge3$、一切满足 \eqref{eq:frame18}~的 $(g,U)$、一切 $\Lam\succeq0$ 且 $\tr\Lam=1$，
\begin{equation}\label{eq:r25main}
2\rho_0-J_w=\Phi-\frac{\delta_0}3 .
\end{equation}
特别地，
\begin{equation}\label{eq:r25equiv}
\kappa_3\le\rho_0\quad\Longleftrightarrow\quad
\Phi\ge\frac{\delta_0}3\ \ \text{对一切 }(g,U,\Lam),
\end{equation}
即 $\Lam_3(G)\le\rho_0n-3$ 对一切 $n$ 与一切图 $G$ 成立
$\iff\Phi\ge\frac{\delta_0}3=\frac{9+7\sqrt5-4\sqrt{35}}{18}=0.054897595005559$。
\end{theorem}

\begin{proof}
由 \eqref{eq:r25ch1}~与 \eqref{eq:r25ch2}，$H=6-2W^2$、$\widetilde E=2-2\|S_w\|_F^2$，故
\[
\frac{c_W}2H+\frac{c_S}2\widetilde E
=3c_W-c_WW^2+c_S-c_S\|S_w\|_F^2 .
\]
代入 \eqref{eq:r25phi}~并用 \eqref{eq:r25key}~把 $|c_Z|Z_4+D_w$ 换成
$c_WW^2+c_S\|S_w\|_F^2-L_w$，$W^2$ 与 $\|S_w\|_F^2$ 的项\textbf{逐项相消}：
\[
\Phi=(1-A_w^2)+(3c_W+c_S)-L_w .
\]
由 \eqref{eq:r25alg}~的第一式 $3c_W+c_S=\varphi$ 与第二式 $\varphi-\frac{\delta_0}3=2\rho_0-1$，
\[
\Phi-\frac{\delta_0}3=1-A_w^2+\varphi-\frac{\delta_0}3-L_w
=2\rho_0-A_w^2-L_w=2\rho_0-J_w .
\]
\eqref{eq:r25equiv}~由 \eqref{eq:r25Jw}~与 $\Phi\ge\delta_0/3$ 的等价性得到：
$\kappa_3\le\rho_0\iff\sup J_w\le2\rho_0\iff J_w\le2\rho_0\ \forall(g,U,\Lam)
\iff\Phi\ge\delta_0/3\ \forall(g,U,\Lam)$。
\end{proof}

\begin{corollary}[判据的尖锐性与冠军处的五项预算]\label{cor:r25sharp}
冠军构型（$W_6=K_1\vee C_5$ 的团吹起帧，$a^*=\frac{3\sqrt{35}+5\sqrt7-7-7\sqrt5}{14}=0.594608575866$、
$z^*=0.500440049768212$）满足 $J_w/2=\rho_0$ 与 $\Phi=\delta_0/3$（差 $6.9\times10^{-17}$），
故 \eqref{eq:r25equiv}~中的常数 $\delta_0/3$ \textbf{不可改进}，判据是尖锐的。
五项在冠军处的取值见表~\ref{tab:r25budget}：五项\emph{全部严格为正}，
其中最大项 $\sin^2\theta$ 只占 $42.63\%$。
\end{corollary}

\begin{table}[ht]\centering\small
\begin{tabular}{llrl}
\toprule
项 & 表达式 & 冠军处取值 & 占 $\delta_0/3$ 的比例\\
\midrule
秩一亏格 & $\sin^2\theta=1-A_w^2=E-\frac{E^2}4$ & $0.023404740945$ & $42.63\%$\\
$h$-色散 & $3c_W\sin^2\theta_1=\frac{c_W}2H$ & $0.010880039083$ & $19.82\%$\\
角向稀释 & $\frac{c_S}2\widetilde E$ & $0.009602398641$ & $17.49\%$\\
四阶 Schoenberg & $|c_Z|Z_4$ & $0.003452901195$ & $6.29\%$\\
主函数亏格 & $D_w$ & $0.007557515142$ & $13.77\%$\\
\midrule
合计 & $\Phi$ & $0.054897595005559$ & $100\%$\\
\bottomrule
\end{tabular}
\caption{冠军处五项亏格的精确预算（$\delta_0/3=0.054897595005559$）。
配套数值：$A_w=0.988228343580$、$E=0.023543312839$、$H=0.087738815$、
$\widetilde E=0.021973300$、$Z_4=0.028299038$、$W=1.719340161947$。}\label{tab:r25budget}
\end{table}

\begin{remark}[这条归约的\textbf{否定性}含义]\label{rem:r25no}
表~\ref{tab:r25budget}~给出一条必须直面的事实：
\emph{任何}在冠军处严格放松的估计都不可能证出 $\Phi\ge\delta_0/3$。
具体地说，
(a) 丢掉任意一项后，剩余四项之和 $\le0.0515<\delta_0/3$（丢 $\sin^2\theta$ 时）——
所以「先证 $\Phi$ 的某一项可以忽略」这条路被封死；
(b) 把 $\widetilde E\le H/3$（引理~\ref{lem:r25chan}\eqref{eq:r25ch3}）用作放缩，
在冠军处 $\widetilde E/H=0.021973/0.087739=0.25044=z^{*2}$，
而 $H/3=0.029246$，故该放缩损失 $29\%$，
把 $3c_W\sin^2\theta_1+\frac{c_S}2\widetilde E$ 换成 $(3c_W+\frac{c_S}6)H$
只能给出 $0.01813<0.02048$，缺口 $2.4\times10^{-3}$；
(c) 用 $q(s)\ge\sqrt s$ 的\emph{任意}放松（例如把 $D_w$ 换成 $0$）损失 $13.77\%$。
这三条把第 22--24 轮那类「矩松弛 $+$ 逐点主函数」的技术路线彻底排除在本问题之外，
与定理~\ref{thm:zonal}（通道型方案的天花板恰为 $\Phi=(3+\sqrt5)/4$）互为印证。
\end{remark}

\subsection{无条件窗口定理}\label{sec:r25:window}

\begin{theorem}[全局窗口归约，无条件]\label{thm:r25window}
令
\begin{equation}\label{eq:r25Wc}
W_c:=\sqrt{3-\frac{\delta_0}{3c_W}}=1.66692753215039379810,\qquad
H_c:=2(3-W_c^2)=\frac{\delta_0}{3c_W}=0.44270520511799510421 .
\end{equation}
若 $W\le W_c$（等价地 $H\ge H_c$、$\sin^2\theta_1\ge\frac{\delta_0}{9c_W}=0.073784200853$、
$\theta_1\ge15.7614316733^\circ$），则 $\Phi\ge\delta_0/3$，从而
$\rho_{g,U}(\Lam)\le\rho_0$。
因此
\begin{equation}\label{eq:r25window}
\kappa_3\le\rho_0\iff
\Bigl[\ \forall(g,U,\Lam)\ \text{满足}\ W>W_c:\ \Phi\ge\frac{\delta_0}3\ \Bigr],
\end{equation}
即全局猜想被压缩到 $\theta_1<15.7614^\circ$（$h$ 接近常值）的薄窗口内。
\end{theorem}

\begin{proof}
丢弃 $\Phi$ 的后四项（均非负），
$\Phi\ge\frac{c_W}2H=c_W(3-W^2)$（用 \eqref{eq:r25ch1}）。
$W\le W_c$ 时 $3-W^2\ge3-W_c^2=\frac{\delta_0}{3c_W}$，故
$\Phi\ge c_W\cdot\frac{\delta_0}{3c_W}=\frac{\delta_0}3$。
\end{proof}

\begin{remark}[窗口的位置与两个端点]\label{rem:r25win}
窗口的两个端点都已被前几轮闭合，这是本轮归约最有价值的结构信息：
\begin{itemize}\itemsep1pt
\item $\theta_1=0$（$h\equiv\frac1{\sqrt3}$，即中心截面 $\Lam=I_3/3$）：
由命题~\ref{prop:r25lat}~与定理~\ref{thm:r24centre}，
$\Phi\ge\frac{\delta_0}3+\frac{m_1}3$，$m_1=1.040\times10^{-3}$，\textbf{带正余量闭合}。
\item $\theta_1\ge15.7614^\circ$：定理~\ref{thm:r25window}，\textbf{无条件闭合}（且只用一项）。
\item 冠军位于窗口\emph{内部}：$\theta_1=6.945557^\circ$、$W=1.719340>W_c$，
$\sin^2\theta_1=0.014623136$，只用了允许上限 $0.073784201$ 的 $19.82\%$。
\end{itemize}
换言之，剩下的问题是：在 $\theta_1\in(0,15.7614^\circ)$ 这条\emph{开区间}上，
证明其余四项之和 $\ge\delta_0/3-3c_W\sin^2\theta_1$。
右端在 $\theta_1\to0$ 时趋于 $\delta_0/3$（此时由第 24 轮已知成立），
在 $\theta_1\to15.7614^\circ$ 时趋于 $0$；冠军处右端为 $0.044017556$，
而其余四项实际给出 $0.031492854$——\textbf{缺口 $1.25\times10^{-2}$}，
这就是本轮之后主方向的准确刻度。
\end{remark}

\subsection{秩一刚性与常纬度归约}\label{sec:r25:rigid}

\begin{proposition}[秩一截面的刚性]\label{prop:r25rank1}
设 $\Lam=zz^\top$（$\|z\|=1$）。则对\textbf{一切}几何 $(g,U)$：
\begin{equation}\label{eq:r25r1}
E=0,\qquad A_w=1,\qquad \sin^2\theta=0,\qquad L_w=\varphi-\Phi ,
\end{equation}
从而
\begin{equation}\label{eq:r25r2}
\kappa_3(\text{秩一})\le\rho_0\iff
L_w\le2\rho_0-1=\varphi-\frac{\delta_0}3=1.56313639374433566 .
\end{equation}
第 26 轮的定理~\ref{thm:r26cone}~将证明锥族上
$\sup L_w=\frac{13+\sqrt5}{10}=\frac{6+\varphi}5=1.523606797750$
（在 $n=5$、$s_0=\frac{\varphi}{4+\varphi}=0.288007155526$ 处取得），
一般帧的数值搜索从未超过此值，故余量为 $3.9530\times10^{-2}$（相对 $2.53\%$）。
\end{proposition}

\begin{proof}
$h_i=\sqrt{u_i^\top zz^\top u_i}=|\langle u_i,z\rangle|=a_i(z)$，故 $E(z)=0$，
$E=\min\le E(z)=0$，而 $E\ge0$，故 $E=0$；引理~\ref{lem:r25dipole}~给 $A_w=1-E/2=1$、
$\sin^2\theta=1-A_w^2=0$。
再由主恒等式 \eqref{eq:r25main}：$2\rho_0-(L_w+1)=\Phi-\frac{\delta_0}3$，
即 $L_w=2\rho_0-1+\frac{\delta_0}3-\Phi=\varphi-\Phi$（用 \eqref{eq:r25alg}~第二、三式）。
\end{proof}

\begin{remark}[秩一截面为何是「最坏」的]\label{rem:r25rank1}
$\sin^2\theta$ 是五项中最大的一项（冠军处占 $42.63\%$），
而命题~\ref{prop:r25rank1}~说明它在秩一截面上\textbf{恒等于零}。
因此秩一截面是 $\Phi$ 的预算中最不利的一面：它必须靠剩下四项凑满 $\delta_0/3$。
好在这四项在秩一截面上并不小：$L_w=\varphi-\Phi$ 说明
$\Phi\ge\delta_0/3\iff L_w\le\varphi-\delta_0/3$，
而 \STW\ 型逐点估计给出的正是 $L_w\le\varphi$（即 $\rho\le\frac{1+\varphi}2=\Phi=1.309017$）。
所以\textbf{秩一截面上需要改进的量恰好是 $\delta_0/3=0.054900$}，
相对 \STW\ 的 $\varphi$ 是 $3.39\%$。这是下一轮 P1 的准确任务。
\end{remark}

\begin{proposition}[常纬度归约到第 24 轮]\label{prop:r25lat}
设 $h\equiv\frac1{\sqrt3}$（即 $\Lam=I_3/3$）。则 $w_i=g_i/\sqrt3$、$W=\sqrt3$、
$\theta_1=0$、$H=\widetilde E=0$、$S_w=I_3/\sqrt3$、$\zeta=\frac1{\sqrt3}$，
且
\begin{equation}\label{eq:r25lat}
\sin^2\theta=\frac12V_w,\qquad
\Phi=\frac12V_w+|c_Z|Z_4+D_w ,
\end{equation}
右端正是第 24 轮的中心截面泛函；由定理~\ref{thm:r24centre}，
$\Phi\ge\frac{\delta_0}3+\frac{m_1}3$，$m_1=1.040\times10^{-3}$。
\end{proposition}

\begin{proof}
$h$ 常值 $\Rightarrow H=\widetilde E=0$；$w=g/\sqrt3\Rightarrow S_w=\frac1{\sqrt3}\sum_ig_iu_iu_i^\top=\frac1{\sqrt3}I_3$，
$W=\sum_ig_i/\sqrt3=3/\sqrt3=\sqrt3$，故 $\cos\theta_1=1$、$\theta_1=0$，$\zeta=z_*^\top S_wz_*=\frac1{\sqrt3}$。
\eqref{eq:r25ch4}~给 $A_w^2=\sqrt3\cdot\frac1{\sqrt3}-\frac12V_w=1-\frac12V_w$，
故 $\sin^2\theta=\frac12V_w$。代入 \eqref{eq:r25phi}~得 \eqref{eq:r25lat}。
最后一步是第 24 轮定理~\ref{thm:r24centre}~的结论（那里记作
$\kappa_3^c\le\rho_0-1.7333\times10^{-4}$，与 $\Phi\ge\frac{\delta_0}3+\frac{m_1}3$ 等价）。
\end{proof}

\begin{remark}[秩一截面即使 $\Lam$ 秩一、$h$ 也不常值]\label{rem:r25lat2}
命题~\ref{prop:r25lat}~与命题~\ref{prop:r25rank1}~是 $\Lam$-单纯形上两个\emph{极端}面：
$\Lam=I_3/3$（谱 $(\frac13,\frac13,\frac13)$）与 $\Lam=zz^\top$（谱 $(1,0,0)$）。
前者已闭合且带余量 $3.47\times10^{-4}$，后者化为一维量 $L_w$ 的上界、余量 $3.94\times10^{-2}$（数值）。
两个面之间的内部（冠军所在的单轴截面 $a^*=0.594609$，谱 $(a^*,b^*,b^*)$）才是难点，
而 \eqref{eq:r25equiv}~说明它等价于同一条 $\Phi\ge\delta_0/3$。
\end{remark}

\subsection{单轴截面的两尺度结构}\label{sec:r25:twoscale}

\begin{proposition}[两尺度精确公式]\label{prop:r25twoscale}
设 $\Lam=bI_3+\rho ee^\top$，$b=\frac{1-\rho}3$，$\rho\in[0,1)$，
$t_i:=\langle u_i,e\rangle$、$s_i:=t_i^2$、$h_i=\sqrt{b+\rho s_i}$。
若 $z=e$ 是 $E(z)$ 的极小点（例如帧关于 $e$ 轴对称时），则
\begin{align}
E&=(1-\rho)^2\sum_ig_i\frac{(\frac13-s_i)^2}{(h_i+|t_i|)^2},\label{eq:r25ts1}\\
H&=\rho^2\sum_{ij}g_ig_j\frac{(s_i-s_j)^2}{(h_i+h_j)^2}.\label{eq:r25ts2}
\end{align}
特别地 $\sin^2\theta=E-\frac{E^2}4$ 含因子 $(1-\rho)^2$，
$3c_W\sin^2\theta_1=\frac{c_W}2H$ 含因子 $\rho^2$，
二者是\textbf{互补的幂次}；且 $E=H=0\iff s_i\equiv\frac13$（常纬度）。
\end{proposition}

\begin{proof}
$h_i-|t_i|=\frac{h_i^2-t_i^2}{h_i+|t_i|}=\frac{b+\rho s_i-s_i}{h_i+|t_i|}
=\frac{(1-\rho)(\frac13-s_i)}{h_i+|t_i|}$（用 $b-(1-\rho)s_i=(1-\rho)(\frac13-s_i)$），
平方加权求和得 \eqref{eq:r25ts1}。
$h_i-h_j=\frac{h_i^2-h_j^2}{h_i+h_j}=\frac{\rho(s_i-s_j)}{h_i+h_j}$，代入 $H$ 的定义得 \eqref{eq:r25ts2}。
$\frac13-s_i$ 与 $s_i-s_j$ 同时为零 $\iff s_i\equiv\frac13$（注意 $\sum_ig_is_i=1=\frac13\sum_ig_i$）。
\end{proof}

\begin{remark}[两尺度之间的张力，本轮的障碍所在]\label{rem:r25tension}
把冠军数（$\rho^*=\frac{3a^*-1}2=0.391912863799$）代入：
\eqref{eq:r25ts1}~给 $E=0.023543313$（与引理~\ref{lem:r25dipole}~直接算出的值一致，残差 $0$），
\eqref{eq:r25ts2}~给 $H=0.087738815$。
按冠军处的预算，含 $h$-色散的三项（$\sin^2\theta$、$3c_W\sin^2\theta_1$、$\frac{c_S}2\widetilde E$）
合计 $0.043887179$，占 $\delta_0/3$ 的 $79.94\%$；
纯角向的两项（$|c_Z|Z_4$、$D_w$）合计 $0.011010416$，占 $20.06\%$。
\textbf{两组各自都不够}（$0.043887<0.054898$、$0.011010<0.054898$），
所以必须同时控制纬度尺度与角向尺度——这正是 $\Phi\ge\delta_0/3$ 的难度来源。
一个有用的精确观察：对\textbf{锥}帧（极点 $+$ 正 $n$ 边形），
只有「极--缘」对满足 $h_i\ne h_j$，而这些对的 $\langle u_i,u_j\rangle^2=z^2=s$ 是常数，故
\begin{equation}\label{eq:r25conetilde}
\widetilde E=s\cdot H\qquad(\text{锥帧，精确}),
\end{equation}
于是锥族上 $\frac{c_W}2H+\frac{c_S}2\widetilde E=(3c_W+3c_Ss)\sin^2\theta_1$，
$h$-色散的两项合并成一个\emph{单参数}系数。冠军处 $s=z^{*2}=0.250440$，
$\widetilde E/H=0.250440$ 与 \eqref{eq:r25conetilde}~完全吻合（数值残差 $<10^{-15}$）。
\end{remark}

\subsection{数值证据}\label{sec:r25:num}

以下均为\textbf{数值}结果，不计入已证结论；全部由
\texttt{证书\_五项亏格恒等式与窗口归约.py} 复现（$48$ 项全绿）。
\begin{enumerate}\itemsep2pt
\item \textbf{锥族}（极点 $+$ 正 $n$ 边形，$n=3,\dots,12$，$240\times240$ 网格 $+$ 冠军精确点）：
$\min\Phi=\delta_0/3$，唯一取等在 $n=5$、$a=a^*$、$z=z^*$（差 $-6.9\times10^{-17}$）。
次小值出现在 $n=8$（$\Phi=0.0979$）与 $n=4$（$\Phi=0.1117$），
即锥族上 $\Phi$ 的第二低点比 $\delta_0/3$ 高 $78\%$。
\item \textbf{一般帧 $\times$ 一般 $\Lam$ 的随机搜索}：$13500$ 个构型
（$N=4,\dots,12$，$\Lam$ 取单轴、随机谱 $+$ 随机正交、随机谱 $+$ 对角三类），
$\min\Phi=0.129025$，余量 $+7.41\times10^{-2}$。
\item \textbf{窗口筛选}：$4000$ 个随机构型中 $791$ 个落在 $W\le W_c$ 内，
其 $\min\Phi=0.212118\gg\delta_0/3$，最大 $W=1.666900\le W_c$，与定理~\ref{thm:r25window}~一致。
\item \textbf{局部极小化}（Powell，变量为 Stiefel 投影 $+$ $\Lam$ 的谱）：
从随机起点出发，$N\ge7$ 时收敛到 $\Phi=0.054897595005560$（与 $\delta_0/3$ 差 $6.7\times10^{-16}$），
且 $\Lam$ 的谱收敛到 $(0.594608582,0.202695712,0.202695707)$，
与冠军单轴谱 $(a^*,b^*,b^*)$ 相符（首元差 $6.0\times10^{-9}$）。
\item \textbf{冠军的严格局部极小性}：$1000$ 个随机扰动（$\varepsilon=3\times10^{-2}$ 到 $10^{-3}$）
无一使 $\Phi$ 低于 $\delta_0/3$，二阶增长率 $(\Phi-\Phi_0)/\varepsilon^2\to1.671$。
\item \textbf{秩一截面}：$500$ 组随机构型给 $L_w\in[1.0139,1.4202]$，
这\emph{只是随机采样的最大值}，不是上确界；沿锥族做一维极大化后
$\sup L_w$ 明显更高（$n=5$ 锥族达 $1.523607$，见第 26 轮定理~\ref{thm:r26cone}），
但仍低于 $2\rho_0-1=1.563136$，余量 $3.9530\times10^{-2}$。
\end{enumerate}

\subsection{本轮状态、诚实标签与方法优先队列（第十二版）}\label{sec:r25:status}

\begin{center}\small
\begin{tabular}{lll}
\toprule
结论 & 状态 & 位置\\
\midrule
坐标函数在 $L^2(\nu)$ 中正交归一、$\|h\|_g=\|a_z\|_g=1$ & \textbf{已证（精确）} & 引理~\ref{lem:r25ortho}\\
偶极子余弦形式 $\sum_iw_ia_i(z)=1-\frac12E(z)$ & \textbf{已证（精确）} & 引理~\ref{lem:r25dipole}\\
\textbf{$A_w=1-\frac E2\le1$}（无条件通道上界） & \textbf{已证} & 引理~\ref{lem:r25dipole}\\
$W^2=3-\frac H2$、$\|S_w\|_F^2=1-\frac{\widetilde E}2$ & \textbf{已证（精确）} & 引理~\ref{lem:r25chan}\\
$\widetilde E\le H/3\iff\|S_w\|_F^2\ge W^2/3$（Schoenberg $Z_2\ge0$） & \textbf{已证} & 引理~\ref{lem:r25chan}\\
$t^4$ 的 Schoenberg 分解与 $M_4$ 公式、$Z_4\ge0$ & \textbf{已证（精确）} & 引理~\ref{lem:r25scho}\\
关键消去 $D_w+|c_Z|Z_4=c_WW^2+c_S\|S_w\|_F^2-L_w$ & \textbf{已证（精确）} & 引理~\ref{lem:r25scho}\\
\textbf{主恒等式 $2\rho_0-J_w=\Phi-\frac{\delta_0}3$} & \textbf{已证（精确）} & 定理~\ref{thm:r25main}\\
$\kappa_3\le\rho_0\iff\Phi\ge\delta_0/3$ & \textbf{已证（等价改写）} & 定理~\ref{thm:r25main}\\
常数 $\delta_0/3$ 不可改进、冠军五项全正 & \textbf{已证（精确取等）} & 推论~\ref{cor:r25sharp}\\
\textbf{窗口定理} $W\le W_c\Rightarrow\rho\le\rho_0$ & \textbf{已证（无条件）} & 定理~\ref{thm:r25window}\\
秩一刚性 $A_w\equiv1$、$L_w=\varphi-\Phi$ & \textbf{已证（精确）} & 命题~\ref{prop:r25rank1}\\
常纬度归约到第 24 轮 & \textbf{已证} & 命题~\ref{prop:r25lat}\\
两尺度公式 \eqref{eq:r25ts1}\eqref{eq:r25ts2} & \textbf{已证（精确）} & 命题~\ref{prop:r25twoscale}\\
锥帧 $\widetilde E=sH$ & \textbf{已证（精确）} & 注~\ref{rem:r25tension}\\
$\sup L_w=1.523744$（秩一截面） & \textbf{数值}（$n=5$ 锥族一维极大化） & \S\ref{sec:r25:num}\\
锥族 $\min\Phi=\delta_0/3$ 唯一取等于冠军 & \textbf{数值}（第 23 轮已解析证明 $\max F=\rho_0$） & \S\ref{sec:r25:num}\\
冠军是 $\Phi$ 的严格局部极小、二阶增长率 $1.671$ & \textbf{数值} & \S\ref{sec:r25:num}\\
\textbf{全局 $\kappa_3\le\rho_0$} & \textbf{未证}（窗口 $\theta_1<15.7614^\circ$ 内） & 注~\ref{rem:r25win}\\
\bottomrule
\end{tabular}
\end{center}

\textbf{方法优先队列（第十二版，数学优先）}：
\begin{enumerate}\itemsep2pt
\item[\textbf{P1}] \textbf{秩一截面的 $L_w\le2\rho_0-1$}（即 $L_w\le\varphi-\delta_0/3$）。
理由：这是五项中最大的一项\emph{恒等于零}的唯一截面，也是 \STW\ 型逐点估计
$L_w\le\varphi$ 与目标之间差距最小、且\textbf{有数值余量 $3.94\times10^{-2}$} 的一个子问题
（对比全局问题余量为零）。入口：命题~\ref{prop:r25rank1}~把问题化为一维量 $L_w$，
而 $L_w$ 又可以精确劈成「符号一致部分 $+$ 负 Gram 质量」（第 26 轮），
于是任务变成一个\emph{单一非负标量}的上界。本项已在第~\ref{sec:r26}~节推进。
\item[\textbf{P2}] \textbf{锥最优性}。定理~\ref{thm:r23uni}~已证明单轴锥族上 $\max F=\rho_0$，
故「极值元必为锥」一旦成立，\textbf{全局猜想立即闭合}（不只是单轴截面）。
注~\ref{rem:r25tension}\eqref{eq:r25conetilde}~说明锥帧上五项中的两项合并为单参数系数，
使锥族成为一维/二维微积分问题。入口：$\mathbf{RP}^2$ 上的 Steiner／Riesz 重排；
极值元一阶条件强制 $n$ 重对称；或用引理~\ref{lem:r25dipole}~的 $z_*$ 作为重排轴。
\item[\textbf{P3}] \textbf{窗口内的两尺度下界}。把
$\sin^2\theta+3c_W\sin^2\theta_1$ 用 \eqref{eq:r25ts1}\eqref{eq:r25ts2}~写成
$(1-\rho)^2\cdot(\text{纬度方差})+\rho^2\cdot(\text{成对纬度方差})$，
对 $s$-分布做 Jensen／LP 对偶。目前的粗估（用 $(h_i+|t_i|)^2\le(1+\sqrt a)^2$、
$(h_i+h_j)^2\le4a$）落在 $\Phi\ge0.0480$，缺口 $6.9\times10^{-3}$（$12.5\%$），
\textbf{已知不足}；改进点是分母的精确处理与 $E/(1-\rho)^2$ 的 LP 对偶。
\item[\textbf{P4}] \textbf{提高中心截面余量}（沿用第十一版 P3）：
$m_1=1.040\times10^{-3}$ 的三处可回收信息中，证书亏格 $D$ 是唯一能带来量级改进的一项。
\item[\textbf{副}] \textbf{下界方向}：把 $\Theta(W_6)=\rho_0$（定理~\ref{thm:r20w6}）
推广成轮族的极值性定理，作为下界一侧的独立结果。
\end{enumerate}

复现：报告目录下 \texttt{证书\_五项亏格恒等式与窗口归约.py}
（$48$ 项全绿，输出 \texttt{\_cert\_r25\_out.txt}，含 $60$ 位常数代数、
四个预备引理的逐条残差、主恒等式在 $800$ 组随机 $(g,U,\Lam)$ 上的核验、
冠军精确取等与 $1000$ 点扰动的局部极小性、窗口定理的代数核对与 $4000$ 点筛选、
秩一刚性、常纬度归约、两尺度公式、锥族 $n=3..10$ 扫描、$13500$ 点随机全局搜索、
局部极小化收敛到冠军谱）；
仓库根目录 \texttt{\_r25b.py}（锥族与随机帧上 $\rho$ 与 $\Phi$ 的对照）、
\texttt{\_r25c.py}（引理逐条残差）、\texttt{\_r25e.py}（冠军核对 $+$ 锥族剖面 $n=3..12$）、
\texttt{\_r25f.py}（随机全局搜索 $+$ Powell 局部极小化）。
\textbf{勘误}：中间草稿 \texttt{\_r25d.py} 的单轴公式与 \texttt{\_r25g.py} 的窗口常数
把 $\sin^2\theta_1=H/6$ 误作 $H/2$，导致 $W_c$ 被算成 $1.710619$（正确值 $1.666928$）、
$\theta_{1,c}$ 被算成 $9.02^\circ$（正确值 $15.7614^\circ$）；
二者均已被合并证书取代，本轮一切结论以 \eqref{eq:r25Wc}~为准。
